\documentclass[10pt,a4paper]{amsart}
\usepackage[margin=19mm]{geometry}
\usepackage{amsmath,amssymb,mathtools}
\usepackage[scr=rsfs,cal=euler]{mathalfa}
\usepackage{xfrac}
\usepackage[unicode,hidelinks,bookmarks=false]{hyperref}
\newcommand{\ie}{i.e.\ }
\newcommand{\cf}{cf.\ }
\newcommand{\resp}{respectively\ }
\newcommand{\Subsection}{\subsection}
\providecommand{\nobreakdash}{\nobreak\hbox{-}}
\setlength{\emergencystretch}{2em}
\multlinegap0pt
\usepackage{bm}
\usepackage{bbm}
\usepackage{dsfont}
\usepackage{xspace}
\usepackage{cancel}
\usepackage{mathtools}
\usepackage{amsthm}
\makeatletter
\@ifundefined{case}{\newtheorem{case}{Case}}{}
\@ifundefined{lemma}{\newtheorem{lemma}{Lemma}}{}
\makeatother
\title[Planar Tur\'an number of two adjacent cycles]{Planar Tur\'an number of two adjacent cycles}
\author{Xinzhe Song, Guiying Yan, Qiang Zhou}
\date{Source: arXiv:2411.18487v2 (CC BY 4.0)}
\begin{document}
\maketitle

\noindent\emph{Source and licence.} Planar Tur\'an number of two adjacent cycles. Version arXiv:2411.18487v2 (CC BY 4.0), distributed under the Creative Commons Attribution 4.0 International licence, as recorded in the arXiv licence metadata for this exact version and shown on the abstract page https://arxiv.org/abs/2411.18487v2. Full rights notice: licenses/arxiv-2411.18487-CC-BY-4.0.txt.\par\smallskip

\noindent\textit{Editorial scope.} Counted unit: the Lemma whose source label is \texttt{face-block} in the article (source0.tex lines 235--242, source label \texttt{face-block}), delivered together with the article's own complete original proof of it (source0.tex lines 244--319, one \texttt{proof} environment of 76 lines). No mathematics is written, repaired or re-derived: statement and proof are the article's own text line for line. Required context retained uncounted: R-v (lines 187--189); partition (lines 197--205). Every definition, display and label of those lines is retained unchanged. Omitted from this package and not redistributed: 1--186, 190--196, 206--234, 243--243, 320--586. Nothing inside a retained excerpt is omitted or altered: every retained body is delivered exactly as the article's lines are written, so no omission receipt of any kind accompanies this package. Citations: the retained excerpts carry no citation command at all, so no bibliography entry of the article is transcribed and no reference is redistributed. Reference adaptation: none was needed and none was made; every reference inside the delivered unit resolves inside it, so no unresolved reference or citation is produced. Printed slips kept literal: no printed word, symbol, hypothesis or number of the counted statement or of its proof was altered. Layout and class: the article's own class is \texttt{article} with packages including geometry, float, latexsym, bm, amsmath, amsfonts, amsthm, amssymb, bbding, mathrsfs, graphicx, hyperref, xcolor; the wrapper is the lane's pinned \texttt{amsart} editorial wrapper at 10pt on A4, which restates the theorem-like environments the retained text needs and adds the article's own macro definitions that the retained text needs (none), with extra wrapper packages (bm, bbm, dsfont, xspace, cancel, mathtools). The delivered numbering is the unit's own. Assets: no retained span contains a figure environment, a graphics-inclusion command, a PDF-page inclusion, a table or a TikZ picture; no image file is copied into this package, the package declares no asset dependency and no image file is redistributed with it. Licence: version arXiv:2411.18487v2 of this article is distributed under the Creative Commons Attribution 4.0 International licence, as recorded in the arXiv licence metadata for this exact version and shown on the abstract page https://arxiv.org/abs/2411.18487v2. A full rights notice is delivered at licenses/arxiv-2411.18487-CC-BY-4.0.txt. Characters: every retained excerpt is pure ASCII, so the delivered engine, XeLaTeX with its default Latin Modern text font, reports no missing character.
\begin{lemma}\label{R-v}
    Given a vertex $v$ in a $(C_3\text{-}C_3)$-free plane graph, if $R_u\subseteq R_v$ for each vertex $u\in V(R_v)$, then $R_v$ is a $3$-face-block and $\sum_{u \in R_v}|\overline{R_u}| \leq 3|V(R_v)|-3$.
\end{lemma}

\begin{lemma}\label{partition}
    Given a vertex $v$ in a $(C_3\text{-}C_3)$-free plane graph $G$ with $R_v\neq\emptyset$, suppose that $(k_1,k_2,...,k_{\ell})$ is a partition of $R_v$. If all the following conditions: 
    \begin{enumerate}
        \item $k_1\leq 2$ when $\ell=3$;
        \item $k_1\leq 3$ and $k_1+k_2\leq 5$ when $\ell=2$;
        \item $k_1\leq 4$ when $\ell=1$;
    \end{enumerate} 
    do not hold, then we have $R_u \subseteq R_v$ for each vertex $u$ in $R_v$.
\end{lemma}

\begin{lemma}\label{face-block}
    For each $3$-face-block $B$ in a $(C_3\textbf{-}C_3)$-free plane graph $G$, we have  
    \begin{enumerate}
        \item $\sum_{v \in B}|\overline{R_v}| \leq 3|V(B)|-3$ when $B\notin\{\mathbf{B}_1,\mathbf{B}_2\}$ or $B\cong\mathbf{B}_3$ and $v_1v_2v_3$ is not a face;
        \item $\sum_{v \in B}|\overline{R_v}| = 3|V(B)|$ when $B\cong\mathbf{B}_2$ or $B\cong\mathbf{B}_1$ and $v_1v_3v_4$ is not a face or $B\cong\mathbf{B}_3$ and $v_1v_2v_3$ is a face;
        \item $\sum_{v \in B}|\overline{R_v}| = 3|V(B)|+3$ when $B\cong\mathbf{B}_1$ and $v_1v_3v_4$ is a face.
    \end{enumerate}
\end{lemma}

\begin{proof}
    It is easy to verify that the statements $2$ and $3$ hold. For each $3$-face-block $B\notin\{\mathbf{B}_1,\mathbf{B}_2,\mathbf{B}_3\}$, we define $|\overline{R_B}|=\mbox{max}\{|\overline{R_u}|:u\in V(B)\}$. From Lemma~\ref{R-v}, if $R_u\subseteq R_v$ for each vertex $u$ of $R_v$, then $R_v=B$ and $\sum_{u \in R_v}|\overline{R_u}| \leq 3|V(R_v)|-3$. 
    Suppose that there is a vertex $u\in V(R_v)$ such that $R_u\not\subseteq R_v$. It follows that there is a $3$-face $F$ containing $u$ such that $F\notin R_v$. From Lemma~\ref{partition}, we have $|\overline{R_B}|\leq 6$. Now we consider the following cases depending on the value of $|\overline{R_B}|$.

    \begin{case} 
        $|\overline{R_B}|=6$.
    \end{case}
    
    From the statement $1$ of Lemma~\ref{partition}, the partition of $R_v$ is $(2,2,2)$. If $v_1=u$, to avoid the appearance of $C_3\textbf{-}C_3$ in $G$, then $F$ must contain $v_5, v_8$ and at least one vertex of $\{v_2,v_3\}$, a contradiction. Similarly, we have $u\notin\{v_1,v_3,v_4,v_6,v_7,v_9\}$. Hence, $u\in\{v_2,v_5,v_8\}$ and $F=v_2v_5v_8$. It implies that $B$ contains at most $(|\overline{R_v}|+1)$ $3$-faces. Thus, $\sum_{v\in B}|\overline{R_v}| \leq 21 < 3\times10-3$. 
    
    \begin{case} 
        $|\overline{R_B}|=5$.
    \end{case}
    
    From the statements $1$ and $2$ of~Lemma \ref{partition}, the partition of $R_v$ is $(2,2,1)$ or $(3,2)$.
    If the former holds, by an argument similar to that above, the $3$-face $F$ is $v_2v_5v_7$ or $v_2v_5v_8$. It follows from the planarity of $G$ that at most one of $v_2v_5v_7$ and $v_2v_5v_8$ is a $3$-face of $G$. Note that $B$ contains at most $(|\overline{R_v}|+1)$ $3$-faces. Thus $\sum_{v \in B}|\overline{R_v}| \leq 18 < 3\times9-3$; 
    If the latter holds, to avoid the appearance of $C_3\textbf{-}C_3$ in $G$, then $v_5v_7, v_1v_3, v_2v_4 \notin E(G)$. Since $G$ is $(C_3\textbf{-}C_3)$-free, $F$ must intersect one vertex of $\{v_5,v_6\}$, $\{v_6,v_7\}$, $\{v_1,v_2\}$ and $\{v_3,v_4\}$, respectively. It follows that $F$ is $v_6v_1v_4$ or $v_6v_2v_3$. If $F$ is $v_6v_1v_4$, then we could find a copy of $C_3\textbf{-}C_3$ which is $F$ and the $3$-face $vv_2v_3$ and an edge connecting them, a contradiction. Then $F$ is $v_6v_2v_3$. By the planarity of $G$, $B$ contains at most $(|\overline{R_v}|+1)$ $3$-faces. Thus, $\sum_{v \in B}|\overline{R_v}| \leq 18 < 3\times8-3$.

    \begin{case} 
        $|\overline{R_B}|=4$.
    \end{case}
    
    From the conditions $1$-$3$ of Lemma~\ref{partition}, the partition of $R_v$ is $(2,1,1)$ or $(2,2)$ or $(3,1)$ or $(4)$.
    
    If the partition of $R_v$ is $(2,1,1)$, by a similar discussion, then $F$ must contain $v_2$, one vertex of $\{v_4,v_5\}$ and one vertex of $\{v_6,v_7\}$, respectively. Thus, $F$ is $v_2v_4v_6$ or $v_2v_4v_7$ or $v_2v_5v_6$ or $v_2v_5v_7$. Since $G$ is planar, $B$ contains at most $(|\overline{R_v}|+1)$ $3$-faces. Then $\sum_{v \in B}|\overline{R_v}| \leq 15 < 3\times8-3$.

    If the partition of $R_v$ is $(3,1)$, to avoid the appearance of $C_3\textbf{-}C_3$ in $G$, then $F$ must contain exact one vertex of $\{v_1,v_2\}$, $\{v_3,v_4\}$ and $\{v_5,v_6\}$, respectively. If $v_1v_3\in E(G)$, then we could find a copy of $C_3\textbf{-}C_3$ which is $F$ and the $3$-face $vv_5v_6$ in $B$ and an edge connecting them, a contradiction. Similarly, we have $v_2v_4, v_1v_4\notin E(G)$. Then $F$ is $v_2v_3v_5$ or $v_2v_3v_6$. But $\{v_2v_3v_5, v_2v_3v_6\}$ cannot exist at the same time by the planarity of $G$. Then $\sum_{v \in B}|\overline{R_v}| \leq 15 < 3\times7-3$.

    If the partition of $R_v$ is $(2,2)$, to avoid the appearance of $C_3\textbf{-}C_3$ in $G$, then $F$ must contain exact one vertex of $\{v_1,v_2\}$, $\{v_2,v_3\}$, $\{v_4,v_5\}$ and $\{v_5,v_6\}$, respectively. By a similar argument, we have $v_1v_3, v_4v_6\notin E(G)$. If $V(F)\subseteq V(R_v)$, then $F$ is $v_2v_4v_5$ or $v_2v_5v_6$ or $v_5v_1v_2$ or $v_5v_2v_3$. Note that there are at most two $3$-faces in $\{v_1v_4v_5, v_1v_4v_6, v_2v_3v_5, v_2v_3v_6\}$ by the planarity of $G$. If $V(F)\not\subseteq V(R_v)$, then $F$ contains $v_2$ and $v_5$. By the maximality of $|\overline{R_B}|$, there are at most two vertices $w_1,w_2\notin V(R_v)$ such that $v_2v_5w_1$ and $v_2v_5w_2$ are $3$-faces. Recall that $G$ is $(C_3\textbf{-}C_3)$-free, $|\overline{R_{w_1}}|=|\overline{R_{w_2}}|=1$(if $w_1$ and $w_2$ exist) and there are at most two choices for $F$ in $\{v_1v_4v_5, v_1v_4v_6, v_2v_3v_5, v_2v_3v_6, v_2v_5w_1, v_2v_5w_2 \}$. Then $\sum_{v \in B}|\overline{R_v}| \leq 18 \leq 3|V(B)|-3$.

    If the partition of $R_v$ is $(4)$, then $R_v$ is a wheel or a fan.
    First, we assume that $R_v$ is a wheel.
    If $F$ contains only one vertex $v_i$ of $R_v$, then we could find a copy of $C_3\textbf{-}C_3$ which is $F$ and the $3$-face $vv_{i+1}v_{i+2}$ and an edge connecting them, where $i\in[4]$ and the subscripts are taken modulo $4$.
    If $F$ contains exact two successive vertices $v_{i}$ and $v_{i+1}$ of $R_v$, then we could find a copy of $C_3\textbf{-}C_3$ which is $F$ and the $3$-face $vv_{i-1}v_{i-2}$ and an edge connecting them, where $i\in[4]$ and the subscripts are taken modulo $4$.
    If $V(F)\subseteq V(R_v)$, then $F$ is $v_1v_2v_3$ or $v_2v_3v_4$ or $v_3v_4v_1$ or $v_4v_1v_2$.
    If $V(F)\not \subseteq V(R_v)$, then $F$ contains $\{v_1, v_3\}$ or $\{v_2, v_4\}$.
    Note that there are at most two vertices $w_1,w_2\notin V(R_v)$ such that $w_1v_1v_{3}$ and $w_2v_1v_3$ or $w_1v_2v_4$ and $w_2v_2v_4$ are $3$-faces.
    By the planarity of $G$, there are at most two choices for $F$ in $\{v_1v_2v_3, v_2v_3v_4, v_3v_4v_1, v_4v_1v_2, w_1v_1v_{3}, w_2v_1v_3, w_1v_2v_4, w_2v_2v_4\}$.
    If both $v_{i}v_{i+1}v_{i+2}$ and $v_{i}v_{i-1}v_{i-2}$ are $3$-faces, then $B\cong\mathbf{B}_1$, a contradiction. 
    If both $w_1v_iv_{i+2}$ and $w_2v_iv_{i+2}$ are $3$-faces, then $|\overline{R_{w_1}}|=|\overline{R_{w_2}}|=1$ by the hypothesis that $G$ is $(C_3\textbf{-}C_3)$-free.
    Thus, $\sum_{v \in B}|\overline{R_v}| = 18 = 3|V(B)|-3$.
    By symmetry, if both $v_{i}v_{i+1}v_{i+2}$ and $w_1v_iv_{i+2}$ are $3$-faces, then $|\overline{R_{w_1}}|=1$ by the hypothesis that $G$ is $(C_3\textbf{-}C_3)$-free.
    Thus, $B\cong\mathbf{B}_2$, a contradiction.
    Now we assume that $R_v$ is a fan. Since $G$ is $(C_3\textbf{-}C_3)$-free, we have $v_1v_3,v_3v_5\notin E(G)$.
    Using a similar argument, we could conclude that $V(F)\subseteq V(R_v)$ or $F$ is the $3$-face $v_2v_4w$, where $w\notin V(R_v)$.
    If $V(F)\subseteq V(R_v)$, then $F$ is $v_2v_3v_4$ or $v_1v_2v_4$ or $v_5v_4v_2$. 
    Note that there are at most two $3$-faces in $\{v_2v_3v_4, v_1v_2v_4, v_5v_4v_2\}$ by the planarity of $G$.
    It follows that there are at most two choices for $F$ in $\{v_2v_3v_4, v_1v_2v_4, v_5v_4v_2, v_2v_4w\}$.
    If both $v_2v_4w$ and $v_2v_3v_4$ are $3$-faces, then $\sum_{v \in B}|\overline{R_v}| \leq 18=3|V(B)|-3$.
    By symmetry of $v_1v_2v_4$ and $v_5v_4v_2$, if both $v_2v_3v_4$ and $v_1v_2v_4$ are $3$-faces, then $B\cong\mathbf{B}_2$, a contradiction. 

    \begin{case}
        $|\overline{R_B}| = 3$. 
    \end{case}

    From the conditions $1$-$3$ of Lemma~\ref{partition}, the partition of $R_v$ is $(1,1,1)$ or $(2,1)$ or $(3)$.
    
    If the partition of $R_v$ is $(1,1,1)$, to avoid the appearance of $C_3\textbf{-}C_3$, then $F$ intersects exact one vertex of $\{v_1,v_2\}$, $\{v_3,v_4\}$ and $\{v_5,v_6\}$, respectively. Note that $B$ contains at most $(|\overline{R_v}|+1)$ $3$-faces. Then $\sum_{v \in B}|\overline{R_v}| \leq 12 < 3|V(B)|-3$.

    If the partition of $R_v$ is $(2,1)$, to avoid the appearance of $C_3\textbf{-}C_3$ in $G$, then $v_1v_3\notin E(G)$ and $F$ intersects at least one vertex of $\{v_1,v_2\}$, $\{v_2,v_3\}$ and $\{v_4,v_5\}$, respectively. 
    We assert that $F$ must contain $v_2$ since $v_1v_3\notin E(G)$, otherwise we could find a copy of $C_3\textbf{-}C_3$ which is the triangle $v_1v_2v_3$ and the $3$-face $vv_4v_5$ and an edge connecting them. 
    Then $F$ must contain $v_2$. By the maximality of $|\overline{R_v}|$ in $B$, there are at most three $3$-faces containing $v_2$. 
    Then $\sum_{v \in B}|\overline{R_v}| \leq 12 < 3|V(B)|-3$.

    If the partition of $R_v$ is $(3)$, then $R_v$ is a wheel or a fan. First, we assume that $R_v$ is a wheel.
    Since $G$ is $(C_3\textbf{-}C_3)$-free, it follows that $F$ intersects at least two vertices of $R_v$.
    If $F$ is $v_1v_2v_3$, then $V(R_v)=V(G)$ and $B\cong\mathbf{B}_3$.
    Hence, $F$ intersects exact two vertices of $R_v$.
    By the maximality of $|\overline{R_B}|$ in $B$, $B$ contains three $3$-faces and four $3$-faces when $|V(B)|=4$ and $|V(B)|=5$, respectively. 
    Then $\sum_{v \in B}|\overline{R_v}| = 3|V(B)|-3$.
    Now we suppose that $R_v$ is a fan. Note that $F$ intersects at least two vertices of $R_v$ by the hypothesis that $G$ is $(C_3\textbf{-}C_3)$-free. 
    If $F$ contains exact two vertices of $R_v$, to avoid the appearance of $C_3\textbf{-}C_3$ in $G$, then $F$ is $wv_1v_3$ or $wv_2v_4$ or $wv_2v_3$, where $w\notin R_v$.
    If $V(F)\subseteq V(R_v)$, then $F$ is $v_1v_2v_3$ or $v_2v_3v_4$ or $v_1v_3v_4$ or $v_1v_2v_4$.
    By the maximality of $|\overline{R_B}|$ in $B$ and the planarity of $G$, $B$ contains at most $(|\overline{R_v}|+1)$ $3$-faces. 
    Then $\sum_{v \in B}|\overline{R_v}| \leq 12 \leq 3|V(B)|-3$.
\end{proof}

\end{document}
